
\documentclass{article}
\usepackage{colortbl}
\usepackage{makecell}
\usepackage{multirow}
\usepackage{supertabular}

\begin{document}

\newcounter{utterance}

\centering \large Interaction Transcript for game `cladder', experiment `full\_v1.5\_default', episode 7123 with reflexion{-}llama{-}3.1{-}8b{-}instruct.
\vspace{24pt}

{ \footnotesize  \setcounter{utterance}{1}
\setlength{\tabcolsep}{0pt}
\begin{supertabular}{c@{$\;$}|p{.15\linewidth}@{}p{.15\linewidth}p{.15\linewidth}p{.15\linewidth}p{.15\linewidth}p{.15\linewidth}}
    \# & \multicolumn{2}{c}{Player} && \multicolumn{2}{c}{Game Master} \\
    \hline

    \theutterance \stepcounter{utterance}  
    & & & \multicolumn{4}{p{0.6\linewidth}}{
        \cellcolor[rgb]{0.9,0.9,0.9}{
            \makecell[{{p{\linewidth}}}]{
                \texttt{\tiny{[P1$\langle$GM]}}
                \texttt{You are an expert in causal inference. The following question is not a typical commonsense query, but rather a meticulously designed question created by a professor specializing in causal inference, intended to assess the students' mastery of the course content.} \\
\texttt{Imagine a self{-}contained, hypothetical world with only the following conditions, and without any unmentioned factors or causal relationships: Jyka has a direct effect on yupt and kwox. Yupt has a direct effect on kwox. Method 1: We look at how jyka correlates with kwox case by case according to yupt. Method 2: We look directly at how jyka correlates with kwox in general. To understand how jyka affects kwox, is it more correct to use the Method 1 than Method 2?} \\
\texttt{Start your answer with 'Yes' or 'No', followed by additional reasoning or evidence to support your explanation.} \\
            }
        }
    }
     \\ \\

    \theutterance \stepcounter{utterance}  
    & \multicolumn{4}{p{0.6\linewidth}}{
        \cellcolor[rgb]{0.561,0.737,0.561}{
            \makecell[{{p{\linewidth}}}]{
                \texttt{\tiny{[P1$\rangle$GM]}}
                \texttt{No.} \\
\\ 
\texttt{In this hypothetical world, we have a situation where there is a confounding variable (yupt) that affects both the treatment (jyka) and the outcome (kwox). This is a classic case of confounding bias, where the observed association between jyka and kwox may be due to the effect of yupt rather than the direct effect of jyka on kwox.} \\
\\ 
\texttt{Method 1, which looks at the correlation between jyka and kwox case by case according to yupt, is actually a form of stratification or matching, which can help to control for the confounding effect of yupt. However, this method still relies on the assumption that yupt is a valid confounder and that the strata are well{-}defined.} \\
\\ 
\texttt{Method 2, which looks directly at the correlation between jyka and kwox in general, is more susceptible to confounding bias because it does not account for the effect of yupt. In this case, the observed association between jyka and kwox may be due to the effect of yupt, rather than the direct effect of jyka on kwox.} \\
\\ 
\texttt{However, the key point here is that yupt also has a direct effect on kwox, which means that yupt is not only a confounder but also an intermediate variable. This is a case of a directed acyclic graph (DAG) with a collider (yupt) and two paths from jyka to kwox: one direct and one indirect through yupt.} \\
\\ 
\texttt{In this case, using Method 1 (stratification by yupt) would actually create a problem known as the "collider problem" or "backdoor path". By conditioning on yupt, we are essentially opening up the backdoor path between jyka and kwox, which can lead to biased estimates of the effect of jyka on kwox.} \\
\\ 
\texttt{Therefore, the correct approach would be to use a method that can block the backdoor path, such as instrumental variables (IV) or a structural equation model (SEM) that takes into account the direct and indirect effects of jyka on kwox through yupt.} \\
            }
        }
    }
    & & \\ \\

    \theutterance \stepcounter{utterance}  
    & & & \multicolumn{2}{p{0.3\linewidth}}{
        \cellcolor[rgb]{0.9,0.9,0.9}{
            \makecell[{{p{\linewidth}}}]{
                \texttt{\tiny{[GM$|$GM]}}
                \texttt{No} \\
            }
        }
    }
    & & \\ \\

    \theutterance \stepcounter{utterance}  
    & & & \multicolumn{2}{p{0.3\linewidth}}{
        \cellcolor[rgb]{0.9,0.9,0.9}{
            \makecell[{{p{\linewidth}}}]{
                \texttt{\tiny{[GM$|$GM]}}
                \texttt{no} \\
            }
        }
    }
    & & \\ \\

    \theutterance \stepcounter{utterance}  
    & & & \multicolumn{2}{p{0.3\linewidth}}{
        \cellcolor[rgb]{0.9,0.9,0.9}{
            \makecell[{{p{\linewidth}}}]{
                \texttt{\tiny{[GM$|$GM]}}
                \texttt{game\_result = WIN} \\
            }
        }
    }
    & & \\ \\

\end{supertabular}
}

\end{document}
